\documentclass[11pt]{article}
\usepackage{amsmath,amssymb}
\usepackage{enumitem}
\usepackage[margin=1in]{geometry}
\usepackage{parskip}

\title{DeLLMa: How It Works}
\author{}
\date{}

\begin{document}
\maketitle

\begin{center}
\fbox{\parbox{0.9\textwidth}{
\small
DeLLMa chooses one action that maximizes \emph{expected utility} under a \emph{fixed} belief over uncertain states. The belief is elicited from an LLM (verbal labels $\to$ scores $\to$ normalize); utilities come from the LLM's rankings of state--action pairs, aggregated via Bradley--Terry. There is no Bayesian update. Two examples: agriculture (crop choice) and stocks (which stock to buy).
}}
\end{center}
\vspace{1em}

%=============================================================================
\section{Notation}
%=============================================================================

\begin{itemize}[leftmargin=*]
\item \textbf{Actions} $A = \{a_1,\ldots,a_n\}$: the options (e.g., plant apples, buy AMD).
\item \textbf{State} $\theta \in \Theta$: a vector of uncertain factors (e.g., climate, demand).
\item \textbf{Context} $C$: all inputs (reports, tables, summaries).
\item \textbf{Utility} $U(\theta,a)$: value of action $a$ in state $\theta$ (higher is better).
\end{itemize}

%=============================================================================
\section{Goal}
%=============================================================================

DeLLMa picks the action that maximizes expected utility:
\begin{equation}\label{eq:EUC}
U_C(a) = \mathbb{E}_{\pi(\theta\mid C)}[U(\theta,a)] = \sum_{\theta \in \Theta} \pi(\theta\mid C)\, U(\theta,a),
\end{equation}
\begin{equation}\label{eq:objective}
a^* = \arg\max_{a \in A} \; U_C(a).
\end{equation}
The distribution $\pi(\theta\mid C)$ is a \textbf{fixed belief} over states given context $C$: it is set once (elicited from the LLM or loaded from cache) and \textbf{not updated} during the pipeline.

%=============================================================================
\section{Mathematical Process}
%=============================================================================

\subsection{State space}

There are $k$ factors (e.g., climate, supply chain, per-product price change). Each factor $f_i$ has $\ell$ possible values (e.g., $\ell=3$). A state is $\theta = (\theta_1,\ldots,\theta_k)$ with $\theta_i \in \{\tilde{f}_i^1,\ldots,\tilde{f}_i^\ell\}$; the full state space has size $\ell^k$ and is not enumerated---states are sampled.

\subsection{Belief $\pi(\theta\mid C)$}

For each factor and each value, the LLM gives a verbal label (e.g., ``very likely'', ``likely''). A fixed score map $V$ converts labels to numbers:

\begin{center}
\begin{tabular}{|l|c|}
\hline
very likely & 6 \\
likely & 5 \\
somewhat likely & 4 \\
somewhat unlikely & 3 \\
unlikely & 2 \\
very unlikely & 1 \\
\hline
\end{tabular}
\end{center}

Scores are normalized per factor, then multiplied across factors (independence):
\begin{equation}\label{eq:pi_i}
\pi_i(\tilde{f}_i^j \mid C) = \frac{V(v_{ij})}{\sum_{j'=1}^{\ell} V(v_{ij'})},
\qquad
\pi(\theta \mid C) = \prod_{i=1}^{k} \pi_i(\theta_i \mid C).
\end{equation}
So the belief is \textbf{elicitation + normalization}; there is no prior/likelihood or Bayesian update.

\subsection{Sampling and state--action pairs}

\begin{itemize}[leftmargin=*]
\item Sample $s$ states from $\pi(\theta\mid C)$ (for each factor, draw from its marginal).
\item Form all pairs $(\theta^{(j)}, a_i)$ for $j=1,\ldots,s$ and each action $a_i$---so each state is paired with every action (variance reduction). Total: $N = s \cdot n$ pairs. Shuffle.
\item Split into overlapping minibatches of size $b$; stride between batch starts is $\lfloor b(1-q) \rfloor$ with overlap $q$.
\end{itemize}

Index the $N$ state--action pairs by $i = 1,\ldots,N$. Let $\text{act}(i)$ be the action in pair $i$, and define
\[
I_a = \{\, i \in \{1,\ldots,N\} : \text{act}(i) = a \,\},
\]
the set of indices of pairs that use action $a$.

\subsection{Rankings to preferences}

The LLM ranks the pairs in each minibatch (e.g., $r_1$ best, $r_b$ worst). Two options:
\begin{itemize}[leftmargin=*]
\item \textbf{Rank2Pairs:} add preference $r_i \succ r_j$ for every $i < j$.
\item \textbf{Top1:} add only $r_1 \succ r_j$ for $j=2,\ldots,b$.
\end{itemize}
Call the set of pairwise comparisons
\[
\Omega \subseteq \{1,\ldots,N\} \times \{1,\ldots,N\},
\]
where $(i,j) \in \Omega$ means ``pair $i$ is preferred to pair $j$.''

\subsection{Bradley--Terry and scores}

Each state--action pair $i$ is treated as an \emph{item} with a (latent) \textbf{strength} or \textbf{utility} parameter $w_i \in \mathbb{R}$. The Bradley--Terry model specifies the probability that item $i$ is preferred to item $j$ as
\begin{equation}\label{eq:BT}
P(i \succ j) = \frac{e^{w_i}}{e^{w_i} + e^{w_j}}.
\end{equation}
Given the observed comparisons $\Omega$, the log-likelihood (with optional $L_2$ regularization) is
\begin{equation}\label{eq:BT_ll}
\mathcal{L}(\mathbf{w})
= \sum_{(i,j) \in \Omega} \log \frac{e^{w_i}}{e^{w_i} + e^{w_j}}
 \;-\; \lambda \sum_{i=1}^{N} w_i^2,
\end{equation}
where $\lambda \ge 0$ controls the regularization strength. The strength vector $\mathbf{w} = (w_1,\ldots,w_N)$ is fitted by maximizing $\mathcal{L}(\mathbf{w})$, using an iterative algorithm (iterative Luce--spectral ranking, ILSR). Optionally, a softmax with temperature $\tau$ is applied to obtain normalized scores $\tilde{w}_i$.

\subsection{Action utility and decision}

For each action $a$, average the (possibly softmaxed) scores over pairs that have that action:
\begin{equation}\label{eq:Ubar}
\bar{U}(a) = \frac{1}{|I_a|} \sum_{i \in I_a} \tilde{w}_i.
\end{equation}
Then
\begin{equation}\label{eq:MC}
a^* = \arg\max_{a \in A} \bar{U}(a).
\end{equation}
This approximates the expected utility in \eqref{eq:EUC} by Monte Carlo over the sampled states; the Bradley--Terry scores play the role of $U(\theta,a)$.

\subsection{Pipeline in one line}

Elicit $\pi(\theta\mid C)$ $\to$ sample states $\to$ form pairs $\to$ minibatches $\to$ LLM ranks $\to$ pairwise $\Omega$ $\to$ ILSR $\to$ $\mathbf{w}$ $\to$ (optional softmax) $\to$ $\bar{U}(a)$ $\to$ $a^* = \arg\max_a \bar{U}(a)$.

%=============================================================================
\section{Data Used}
%=============================================================================

DeLLMa uses only:
\begin{itemize}[leftmargin=*]
\item \textbf{Context $C$:} reports, summaries, tables fed into the LLM.
\item \textbf{Cached beliefs (optional):} a precomputed $\pi$ stored in JSON, so the LLM is not queried again for the belief step.
\end{itemize}
There is no observation that triggers a Bayesian update.

%=============================================================================
\section{Example 1: Agriculture}
%=============================================================================

\textbf{Context:} USDA fruit report and California price/yield statistics for the chosen fruits. \\
\textbf{States:} climate, supply chain, per-fruit price and yield change (discretized). \\
\textbf{Actions:} e.g., ``Plant 10 acres of apples'', ``Plant 10 acres of avocados'', \ldots \\
\textbf{Flow:} $C$ (and optionally cached beliefs) $\to$ $\pi(\theta\mid C)$ $\to$ sample $\to$ rank $\to$ Bradley--Terry $\to$ $\bar{U}(a)$ $\to$ $a^*$.

%=============================================================================
\section{Example 2: Stocks}
%=============================================================================

\textbf{Context:} Historical monthly stock prices and current price for the chosen tickers; buy/sell dates. \\
\textbf{States:} economic health, market sentiment, per-stock factors (growth, launches, etc.). \\
\textbf{Actions:} e.g., ``Buy AMD'', ``Buy GOOGL'', \ldots \\
\textbf{Flow:} Same as above.

%=============================================================================
\section{Extension: Power Grid (Estimating Inertia $M$ in IEEE 14-Bus)}
%=============================================================================

DeLLMa can be extended to \emph{active probing} in power systems, for example to choose a small excitation that best helps estimate the inertia $M$ of generators in the IEEE 14-bus test system.

\textbf{Setup (IEEE 14-bus inertia estimation):}
\begin{itemize}[leftmargin=*]
\item \textbf{State $\theta$:} unknown dynamic parameters and operating conditions, e.g., generator inertias $M$, damping $D$, loading level, and which lines or units are in/out of service.
\item \textbf{Context $C$:} known grid information: IEEE 14-bus topology and nominal parameters, current operating point (generation, loads, voltages), any past disturbance--response data, and a description of the identification method (how $M$ will be estimated from frequency measurements).
\item \textbf{Actions $A$:} candidate \emph{active probes}, such as:
  \begin{itemize}
  \item $a_1$: 1\% step increase in active power setpoint of a chosen generator for 2 seconds;
  \item $a_2$: 2\% pulse on another generator;
  \item $a_3$: small oscillatory modulation of load at a selected bus;
  \item etc.
  \end{itemize}
  Each $a \in A$ is assumed to be safe and within operational limits.
\item \textbf{Utility $U(\theta,a)$:} quality of the inertia estimate after applying probe $a$ in state $\theta$, for example
\[
U(\theta,a) = - \operatorname{Var}\bigl(\widehat{M} \mid \theta, a\bigr),
\]
so higher utility means a more informative probe (lower expected variance of the estimator $\widehat{M}$) given the same safety constraints.
\end{itemize}

\textbf{How DeLLMa applies:} using the general pipeline,
\begin{enumerate}[leftmargin=*]
\item Use $C$ to elicit a fixed belief $\pi(\theta\mid C)$ over the unknown dynamic state (plausible values of $M$, $D$, loading, topology).
\item Sample states $\theta^{(1)},\ldots,\theta^{(s)} \sim \pi(\theta\mid C)$ and form state--action pairs $(\theta^{(j)},a_i)$ for all $j$ and $a_i \in A$.
\item For each pair, describe to the LLM the operating condition and the probe (e.g., ``IEEE 14-bus near base load, inertia likely in [X,Y]; apply probe $a_i$.''), and let the LLM rank the pairs in minibatches according to how well they will allow $M$ to be estimated.
\item Convert rankings to pairwise comparisons $\Omega$, fit Bradley--Terry to obtain strengths $w_i$ for all pairs, and compute $\bar{U}(a)$ via \eqref{eq:Ubar}.
\item Choose the probe
\[
a^* = \arg\max_{a \in A} \bar{U}(a),
\]
that is, the active probe expected (under $\pi(\theta\mid C)$) to give the most informative frequency response for estimating inertia $M$ in the IEEE 14-bus system.
\end{enumerate}

\textbf{Caveats.} (1) Probes must satisfy all operational constraints (voltage, line flows, frequency); only safe actions should be included in $A$, and they may be further screened by a power-flow or dynamic simulation. (2) The procedure is suited to \emph{planning} experiments (e.g., off-line studies or scheduled tests), not real-time emergency control. (3) Quantitative evaluation requires a simulator or historical disturbance data to link the chosen probe to actual estimation accuracy for $M$.

%=============================================================================
\section{DeLLMa vs.\ Bayesian Optimization}
%=============================================================================

\begin{center}
\begin{tabular}{|l|c|c|}
\hline
& \textbf{DeLLMa} & \textbf{Bayesian optimization} \\
\hline
Goal & One action $a^*$ (max expected utility) & Sequence of points (max unknown $f$) \\
\hline
Belief & Over states $\theta$; fixed, not updated & Over objective $f$; prior $\to$ posterior \\
\hline
Update & None & After each observation $(x_i,y_i)$ \\
\hline
\end{tabular}
\end{center}

DeLLMa is \textbf{not} Bayesian optimization: it uses a fixed belief and makes one decision; BO updates a prior to a posterior and designs a sequence of experiments.

%=============================================================================
\section{Summary}
%=============================================================================

DeLLMa maximizes expected utility under a \textbf{fixed} belief $\pi(\theta\mid C)$ elicited from an LLM. Utilities are inferred from the LLM's rankings via Bradley--Terry. No Bayesian update; the only ``data'' is the context $C$ (and optionally cached beliefs). It is a decision-under-uncertainty procedure, not sequential optimization like BO.

\vspace{0.5em}
\noindent\textit{Liu et al. (2024), DeLLMa. arXiv:2402.02392.}

\end{document}
